
\begin{center}{\LARGE\bfseries\color{ink} Automatismes -- Série 1}\\[2pt]{\small 20 minutes -- sans calculatrice -- une seule réponse exacte par question}\end{center}
\begin{enumerate}[label=\textbf{\arabic*.},itemsep=6pt,leftmargin=*]
\item La moitié de l'inverse de $4$ est égale à :\\[2pt]\begin{tabular}{@{}*{4}{p{4.30cm}}}(A) $\frac12$ & (B) $\frac18$ & (C) $2$ & (D) $8$\end{tabular}
\item On considère $F=a+\dfrac{b}{cd}$. Pour $a=\frac32,\ b=-1,\ c=2,\ d=\frac14$, la valeur de $F$ est :\\[2pt]\begin{tabular}{@{}*{4}{p{4.30cm}}}(A) $\frac12$ & (B) $\frac72$ & (C) $-\frac72$ & (D) $-\frac12$\end{tabular}
\item Le prix d'un article baisse de $20\,\%$ puis augmente de $20\,\%$. Au total, ce prix a connu :\\[2pt]\begin{tabular}{@{}*{2}{p{8.60cm}}}(A) une baisse de $4\,\%$ & (B) aucune évolution\\[2pt] (C) une hausse de $4\,\%$ & (D) une baisse de $40\,\%$\end{tabular}
\item On lance un dé à 4 faces. Les probabilités des faces 1, 2, 3 sont $0{,}15$, $\frac15$, $0{,}4$ ; celle de la face 4 est $x$. Alors :\\[2pt]\begin{tabular}{@{}*{4}{p{4.30cm}}}(A) $x=0{,}35$ & (B) $x=\frac34$ & (C) $x=\frac14$ & (D) $x=0{,}6$\end{tabular}
\item La forme développée et réduite de $(x+4)^2-(3x-2)^2$ est :\\[2pt]\begin{tabular}{@{}*{2}{p{8.60cm}}}(A) $-8x^2-4x+20$ & (B) $-8x^2+20x+12$\\[2pt] (C) $10x^2+20x+12$ & (D) $-8x^2+20x+20$\end{tabular}
\item On sait que $P(\overline A)=0{,}35$, $P(B)=0{,}5$ et $P(A\cap B)=0{,}2$. La probabilité $P(A\cup B)$ est égale à :\\[2pt]\begin{tabular}{@{}*{4}{p{4.30cm}}}(A) $0{,}95$ & (B) $1{,}15$ & (C) $0{,}85$ & (D) $0{,}7$\end{tabular}
\item \begin{minipage}[t]{0.50\linewidth}On considère l'arbre ci-contre. La probabilité de $A$ sachant $B$, notée $P_B(A)$, est égale à :\\[4pt]\begin{tabular}{@{}*{2}{p{4.15cm}}}(A) $0{,}32$ & (B) $0{,}6$\\[2pt] (C) $\frac9{16}$ & (D) $0{,}18$\end{tabular}\end{minipage}\hfill\begin{minipage}[t]{0.46\linewidth}\vspace{-8pt}\arbre{$0{,}3$}{$0{,}6$}{$0{,}4$}{$0{,}7$}{$0{,}2$}{$0{,}8$}\end{minipage}
\item La forme canonique de $2x^2-12x+5$ est :\\[2pt]\begin{tabular}{@{}*{2}{p{8.60cm}}}(A) $2(x+3)^2-13$ & (B) $2(x-3)^2-13$\\[2pt] (C) $2(x-3)^2+5$ & (D) $2(x-3)^2-4$\end{tabular}
\item Les racines de $2x^2+3x-2$ sont :\\[2pt]\begin{tabular}{@{}*{4}{p{4.30cm}}}(A) $\{-\frac12\,;2\}$ & (B) $\{\frac12\,;2\}$ & (C) $\{-\frac12\,;-2\}$ & (D) $\{\frac12\,;-2\}$\end{tabular}
\item La fonction $f(x)=-(x+1)^2+4$ admet :\\[2pt]\begin{tabular}{@{}*{2}{p{8.60cm}}}(A) un maximum égal à $4$, atteint en $x=-1$ & (B) un maximum égal à $4$, atteint en $x=1$\\[2pt] (C) un minimum égal à $4$, atteint en $x=-1$ & (D) un maximum égal à $-1$, atteint en $x=4$\end{tabular}
\end{enumerate}
\newpage
\begin{center}{\LARGE\bfseries\color{ink} Automatismes -- Série 2}\\[2pt]{\small 20 minutes -- sans calculatrice -- une seule réponse exacte par question}\end{center}
\begin{enumerate}[label=\textbf{\arabic*.},itemsep=6pt,leftmargin=*]
\item L'inverse de la moitié de $\frac13$ est égal à :\\[2pt]\begin{tabular}{@{}*{4}{p{4.30cm}}}(A) $\frac16$ & (B) $\frac23$ & (C) $6$ & (D) $\frac32$\end{tabular}
\item On considère $G=\dfrac{a+b}{c-d}$. Pour $a=\frac12,\ b=\frac14,\ c=\frac13,\ d=-\frac16$, la valeur de $G$ est :\\[2pt]\begin{tabular}{@{}*{4}{p{4.30cm}}}(A) $\frac32$ & (B) $\frac92$ & (C) $\frac94$ & (D) $\frac12$\end{tabular}
\item Un article passe de $80$ \texteuro{} à $92$ \texteuro. Son prix a augmenté de :\\[2pt]\begin{tabular}{@{}*{4}{p{4.30cm}}}(A) $12\,\%$ & (B) $15\,\%$ & (C) $11{,}5\,\%$ & (D) $8\,\%$\end{tabular}
\item Les faces 1 et 2 d'un dé à 4 faces ont pour probabilités $\frac12$ et $\frac16$ ; les faces 3 et 4 ont la même probabilité $x$. Alors :\\[2pt]\begin{tabular}{@{}*{4}{p{4.30cm}}}(A) $x=\frac13$ & (B) $x=\frac1{12}$ & (C) $x=\frac14$ & (D) $x=\frac16$\end{tabular}
\item La forme développée et réduite de $(2x+3)^2-(x-1)^2$ est :\\[2pt]\begin{tabular}{@{}*{2}{p{8.60cm}}}(A) $3x^2+10x+8$ & (B) $3x^2+14x+10$\\[2pt] (C) $3x^2+14x+8$ & (D) $5x^2+10x+8$\end{tabular}
\item On sait que $P(A)=0{,}7$, $P(B)=0{,}6$ et $P(A\cup B)=0{,}9$. La probabilité $P(A\cap B)$ est égale à :\\[2pt]\begin{tabular}{@{}*{4}{p{4.30cm}}}(A) $0{,}1$ & (B) $0{,}4$ & (C) $1{,}3$ & (D) $0{,}42$\end{tabular}
\item \begin{minipage}[t]{0.50\linewidth}On considère l'arbre ci-contre. La probabilité $P(B)$ est égale à :\\[4pt]\begin{tabular}{@{}*{2}{p{4.15cm}}}(A) $0{,}6$ & (B) $0{,}4$\\[2pt] (C) $0{,}2$ & (D) $0{,}3$\end{tabular}\end{minipage}\hfill\begin{minipage}[t]{0.46\linewidth}\vspace{-8pt}\arbre{$0{,}5$}{$0{,}4$}{$0{,}6$}{$0{,}5$}{$0{,}2$}{$0{,}8$}\end{minipage}
\item La forme canonique de $-3x^2+6x+1$ est :\\[2pt]\begin{tabular}{@{}*{2}{p{8.60cm}}}(A) $-3(x+1)^2+4$ & (B) $-3(x-1)^2-2$\\[2pt] (C) $-3(x-1)^2+4$ & (D) $3(x-1)^2+4$\end{tabular}
\item Les racines de $3x^2-7x+2$ sont :\\[2pt]\begin{tabular}{@{}*{4}{p{4.30cm}}}(A) $\{2\,;\frac13\}$ & (B) $\{-2\,;-\frac13\}$ & (C) $\{2\,;-\frac13\}$ & (D) $\{1\,;\frac23\}$\end{tabular}
\item La fonction $f(x)=x^2-4x+7$ admet :\\[2pt]\begin{tabular}{@{}*{2}{p{8.60cm}}}(A) un minimum égal à $7$ & (B) un minimum égal à $3$\\[2pt] (C) un maximum égal à $3$ & (D) un minimum égal à $2$\end{tabular}
\end{enumerate}
\newpage
\begin{center}{\LARGE\bfseries\color{ink} Automatismes -- Série 3}\\[2pt]{\small 20 minutes -- sans calculatrice -- une seule réponse exacte par question}\end{center}
\begin{enumerate}[label=\textbf{\arabic*.},itemsep=6pt,leftmargin=*]
\item L'opposé du carré de l'inverse de $2$ est égal à :\\[2pt]\begin{tabular}{@{}*{4}{p{4.30cm}}}(A) $\frac14$ & (B) $-\frac12$ & (C) $-4$ & (D) $-\frac14$\end{tabular}
\item On considère $F=a+\dfrac{b}{cd}$. Pour $a=-\frac13,\ b=2,\ c=-3,\ d=\frac12$, la valeur de $F$ est :\\[2pt]\begin{tabular}{@{}*{4}{p{4.30cm}}}(A) $1$ & (B) $-1$ & (C) $-\frac53$ & (D) $\frac53$\end{tabular}
\item Le prix d'un article augmente de $10\,\%$, puis de nouveau de $10\,\%$. Au total, il a augmenté de :\\[2pt]\begin{tabular}{@{}*{4}{p{4.30cm}}}(A) $21\,\%$ & (B) $20\,\%$ & (C) $11\,\%$ & (D) $10{,}1\,\%$\end{tabular}
\item Les probabilités des faces 1, 2, 3 d'un dé à 4 faces sont $\frac13$, $0{,}1$ et $\frac15$ ; celle de la face 4 est $x$. Alors :\\[2pt]\begin{tabular}{@{}*{4}{p{4.30cm}}}(A) $x=\frac{19}{30}$ & (B) $x=\frac{11}{30}$ & (C) $x=0{,}7$ & (D) $x=\frac7{30}$\end{tabular}
\item La forme développée et réduite de $(x-5)^2-(2x-1)^2$ est :\\[2pt]\begin{tabular}{@{}*{2}{p{8.60cm}}}(A) $-3x^2-14x+24$ & (B) $-3x^2-6x+26$\\[2pt] (C) $-3x^2+6x+24$ & (D) $-3x^2-6x+24$\end{tabular}
\item Les événements $A$ et $B$ sont indépendants, avec $P(A)=0{,}6$ et $P(B)=0{,}5$. La probabilité $P(A\cup B)$ est égale à :\\[2pt]\begin{tabular}{@{}*{4}{p{4.30cm}}}(A) $0{,}3$ & (B) $1{,}1$ & (C) $0{,}8$ & (D) $0{,}9$\end{tabular}
\item \begin{minipage}[t]{0.50\linewidth}On considère l'arbre ci-contre. La probabilité $P_B(A)$ est égale à :\\[4pt]\begin{tabular}{@{}*{2}{p{4.15cm}}}(A) $\frac13$ & (B) $0{,}3$\\[2pt] (C) $0{,}1$ & (D) $0{,}5$\end{tabular}\end{minipage}\hfill\begin{minipage}[t]{0.46\linewidth}\vspace{-8pt}\arbre{$0{,}2$}{$0{,}5$}{$0{,}5$}{$0{,}8$}{$0{,}25$}{$0{,}75$}\end{minipage}
\item La forme canonique de $x^2-3x+1$ est :\\[2pt]\begin{tabular}{@{}*{2}{p{8.60cm}}}(A) $(x+\frac32)^2-\frac54$ & (B) $(x-\frac32)^2+1$\\[2pt] (C) $(x-3)^2-8$ & (D) $(x-\frac32)^2-\frac54$\end{tabular}
\item L'ensemble des solutions de l'inéquation $x^2-x-2>0$ est :\\[2pt]\begin{tabular}{@{}*{2}{p{8.60cm}}}(A) $]-1\,;2[$ & (B) $]-\infty\,;-1[\,\cup\,]2\,;+\infty[$\\[2pt] (C) $]-\infty\,;-2[\,\cup\,]1\,;+\infty[$ & (D) $]-2\,;1[$\end{tabular}
\item Le sommet de la parabole représentant $f(x)=2x^2-8x+3$ est le point :\\[2pt]\begin{tabular}{@{}*{4}{p{4.30cm}}}(A) $S(-2\,;-5)$ & (B) $S(2\,;3)$ & (C) $S(2\,;-5)$ & (D) $S(-2\,;27)$\end{tabular}
\end{enumerate}
\ifdefined\eleve\else
\newpage
\begin{center}{\LARGE\bfseries\color{warn} Corrigé}\end{center}
\small
\textbf{Série 1.}\quad 1.\ (B), 2.\ (D), 3.\ (A), 4.\ (C), 5.\ (B), 6.\ (A), 7.\ (C), 8.\ (B), 9.\ (D), 10.\ (A)
\begin{enumerate}[label=\textbf{\arabic*.},itemsep=1pt,leftmargin=*,topsep=2pt]
\item $\Rightarrow$ $\frac18$ \quad {\color{gray}inverse de $4$ : $\frac14$ ; moitié : $\frac18$.}
\item $\Rightarrow$ $-\frac12$ \quad {\color{gray}$cd=\frac12$, $\frac{-1}{1/2}=-2$, $\frac32-2=-\frac12$.}
\item $\Rightarrow$ une baisse de $4\,\%$ \quad {\color{gray}$0{,}8\times1{,}2=0{,}96$.}
\item $\Rightarrow$ $x=\frac14$ \quad {\color{gray}$0{,}15+0{,}2+0{,}4=0{,}75$, donc $x=\frac14$.}
\item $\Rightarrow$ $-8x^2+20x+12$ \quad {\color{gray}$x^2+8x+16-9x^2+12x-4$.}
\item $\Rightarrow$ $0{,}95$ \quad {\color{gray}$P(A)=0{,}65$ ; $0{,}65+0{,}5-0{,}2=0{,}95$.}
\item $\Rightarrow$ $\frac9{16}$ \quad {\color{gray}$P(A\cap B)=0{,}18$, $P(\overline A\cap B)=0{,}14$, $P(B)=0{,}32$ ; $\frac{0{,}18}{0{,}32}=\frac9{16}$.}
\item $\Rightarrow$ $2(x-3)^2-13$ \quad {\color{gray}$\alpha=3$, $\beta=18-36+5=-13$.}
\item $\Rightarrow$ $\{\frac12\,;-2\}$ \quad {\color{gray}$\Delta=25$, $x=\frac{-3\pm5}{4}$.}
\item $\Rightarrow$ un maximum égal à $4$, atteint en $x=-1$ \quad {\color{gray}$a<0$ : maximum ; sommet $(-1\,;4)$.}
\end{enumerate}\medskip
\textbf{Série 2.}\quad 1.\ (C), 2.\ (A), 3.\ (B), 4.\ (D), 5.\ (C), 6.\ (B), 7.\ (D), 8.\ (C), 9.\ (A), 10.\ (B)
\begin{enumerate}[label=\textbf{\arabic*.},itemsep=1pt,leftmargin=*,topsep=2pt]
\item $\Rightarrow$ $6$ \quad {\color{gray}moitié de $\frac13$ : $\frac16$ ; inverse : $6$.}
\item $\Rightarrow$ $\frac32$ \quad {\color{gray}$a+b=\frac34$, $c-d=\frac12$, $\frac34\div\frac12=\frac32$.}
\item $\Rightarrow$ $15\,\%$ \quad {\color{gray}$\frac{92-80}{80}=\frac{12}{80}=0{,}15$.}
\item $\Rightarrow$ $x=\frac16$ \quad {\color{gray}$2x=1-\frac12-\frac16=\frac13$.}
\item $\Rightarrow$ $3x^2+14x+8$ \quad {\color{gray}$4x^2+12x+9-x^2+2x-1$.}
\item $\Rightarrow$ $0{,}4$ \quad {\color{gray}$P(A\cap B)=0{,}7+0{,}6-0{,}9$.}
\item $\Rightarrow$ $0{,}3$ \quad {\color{gray}$0{,}5\times0{,}4+0{,}5\times0{,}2=0{,}2+0{,}1$.}
\item $\Rightarrow$ $-3(x-1)^2+4$ \quad {\color{gray}$\alpha=1$, $\beta=-3+6+1=4$.}
\item $\Rightarrow$ $\{2\,;\frac13\}$ \quad {\color{gray}$\Delta=25$, $x=\frac{7\pm5}{6}$.}
\item $\Rightarrow$ un minimum égal à $3$ \quad {\color{gray}$(x-2)^2+3$ : minimum $3$.}
\end{enumerate}\medskip
\textbf{Série 3.}\quad 1.\ (D), 2.\ (C), 3.\ (A), 4.\ (B), 5.\ (D), 6.\ (C), 7.\ (A), 8.\ (D), 9.\ (B), 10.\ (C)
\begin{enumerate}[label=\textbf{\arabic*.},itemsep=1pt,leftmargin=*,topsep=2pt]
\item $\Rightarrow$ $-\frac14$ \quad {\color{gray}inverse de $2$ : $\frac12$ ; carré : $\frac14$ ; opposé : $-\frac14$.}
\item $\Rightarrow$ $-\frac53$ \quad {\color{gray}$cd=-\frac32$, $\frac{2}{-3/2}=-\frac43$, $-\frac13-\frac43$.}
\item $\Rightarrow$ $21\,\%$ \quad {\color{gray}$1{,}1^2=1{,}21$.}
\item $\Rightarrow$ $x=\frac{11}{30}$ \quad {\color{gray}$0{,}1+\frac15=\frac3{10}$ ; $\frac13+\frac3{10}=\frac{19}{30}$ ; $x=\frac{11}{30}$.}
\item $\Rightarrow$ $-3x^2-6x+24$ \quad {\color{gray}$x^2-10x+25-4x^2+4x-1$.}
\item $\Rightarrow$ $0{,}8$ \quad {\color{gray}$P(A\cap B)=0{,}3$ ; $0{,}6+0{,}5-0{,}3$.}
\item $\Rightarrow$ $\frac13$ \quad {\color{gray}$P(A\cap B)=0{,}1$, $P(\overline A\cap B)=0{,}2$, $P(B)=0{,}3$ ; $\frac{0{,}1}{0{,}3}=\frac13$.}
\item $\Rightarrow$ $(x-\frac32)^2-\frac54$ \quad {\color{gray}$\alpha=\frac32$, $\beta=\frac94-\frac92+1=-\frac54$.}
\item $\Rightarrow$ $]-\infty\,;-1[\,\cup\,]2\,;+\infty[$ \quad {\color{gray}$x^2-x-2=(x+1)(x-2)$, positif à l'extérieur des racines.}
\item $\Rightarrow$ $S(2\,;-5)$ \quad {\color{gray}$\alpha=2$, $\beta=8-16+3=-5$.}
\end{enumerate}\medskip
\fi